\chapter{Project Management}
\label{chap:project_management}

\section{Project Scope}

Project scope defines the boundaries of the Horse Racing Tournament Management System and identifies which functionalities are included in the project and which are intentionally excluded.

The project focuses on developing a web-based system that supports the management of horse racing tournaments. The system covers user management, horse management, jockey management, tournament management, race scheduling, result processing, ranking calculation, and spectator viewing functions.

Several advanced functionalities such as online payment processing, live video streaming, AI-based prediction systems, and mobile applications are excluded from the current project scope due to time and resource constraints.

Clearly defining the project scope helps the development team maintain focus and avoid scope creep during implementation.

\section{Minimum Viable Product (MVP)}

The MVP represents the minimum set of features required to demonstrate the complete business workflow of a horse racing tournament.

The MVP includes:

\begin{itemize}
	\item User authentication and authorization
	\item Horse registration and management
	\item Jockey invitation and confirmation
	\item Tournament creation and management
	\item Race scheduling
	\item Result management
	\item Ranking calculation
	\item Spectator viewing functions
\end{itemize}

The MVP workflow begins with tournament creation, followed by horse registration, jockey assignment, race scheduling, result confirmation, and ranking publication.

The MVP ensures that all major stakeholders can successfully interact with the system during the demonstration phase.

\section{Team Organization}

The project is developed by Group 3 consisting of seven members.

The team structure follows a collaborative approach where responsibilities are distributed among members according to their technical skills and project requirements.

The Team Leader is responsible for project planning, task assignment, progress monitoring, and communication with the instructor.

Other members contribute to requirement analysis, system design, frontend development, backend development, database design, testing, and documentation activities.

A clear responsibility structure helps improve collaboration efficiency and ensures that project deliverables are completed on time.

\begin{table}[H]
	\centering
	\caption{Team Responsibilities}
	\begin{tabular}{|c|p{4cm}|p{7cm}|}
		\hline
		\textbf{No.} & \textbf{Role} & \textbf{Responsibility} \\
		\hline
		1 & Team Leader & Project planning, coordination and monitoring \\
		\hline
		2 & Member & Requirement analysis and documentation \\
		\hline
		3 & Member & UML diagrams and system design \\
		\hline
		4 & Member & Frontend development \\
		\hline
		5 & Member & Backend development \\
		\hline
		6 & Member & Database design and implementation \\
		\hline
		7 & Member & Testing and quality assurance \\
		\hline
	\end{tabular}
\end{table}

\section{Project Planning}

The project development process is divided into several phases.

\begin{enumerate}
	\item Requirement Analysis
	\item System Design
	\item Database Design
	\item User Interface Design
	\item Backend Development
	\item Frontend Development
	\item Testing and Integration
	\item Documentation and Presentation
\end{enumerate}

Each phase contains specific deliverables and milestones that are tracked throughout the project lifecycle.

Project planning enables the team to allocate resources effectively and monitor project progress continuously.

\begin{table}[H]
	\centering
	\caption{Project Development Phases}
	\begin{tabular}{|c|p{4cm}|p{7cm}|}
		\hline
		\textbf{Phase} & \textbf{Objective} & \textbf{Expected Output} \\
		\hline
		1 & Requirement Analysis & Scope, MVP, Use Cases \\
		\hline
		2 & System Design & UML Diagrams, Architecture Design \\
		\hline
		3 & Database Design & ERD and Database Schema \\
		\hline
		4 & Development & Frontend and Backend Modules \\
		\hline
		5 & Testing & Test Cases and Bug Fixes \\
		\hline
		6 & Documentation & Final Report and Presentation \\
		\hline
	\end{tabular}
\end{table}

\section{Agile Development Process}

The project adopts the Agile development methodology.

Agile promotes iterative development, continuous feedback, and flexibility when handling project changes.

The development process is organized into small tasks that can be completed independently and reviewed regularly.

Regular meetings are conducted to discuss project progress, identify issues, and coordinate upcoming activities.

This approach allows the team to adapt quickly to changing requirements while maintaining development efficiency.

\section{Task Management with Jira}

Jira is used as the primary project management tool throughout the project lifecycle.

Project functionalities are divided into Epics based on system modules, including:

\begin{itemize}
	\item Authentication and Authorization
	\item Horse Management
	\item Jockey Management
	\item Tournament Management
	\item Registration Management
	\item Race Scheduling
	\item Result and Ranking Management
	\item Spectator View
\end{itemize}

Each Epic contains multiple tasks assigned to specific team members.

The Kanban workflow is used to monitor task status and project progress.

\begin{figure}[H]
	\centering
	% \includegraphics[width=\textwidth]{figures/jira_board.png}
	\caption{Jira Kanban Board}
	\label{fig:jira_board}
\end{figure}

The task workflow includes the following stages:

\begin{itemize}
	\item To Do
	\item In Progress
	\item Review
	\item Done
\end{itemize}

Using Jira improves transparency, task tracking, and collaboration among team members.

\section{Risk Management}

Several project risks were identified during the planning phase.

\begin{table}[H]
	\centering
	\caption{Project Risk Assessment}
	\begin{tabular}{|p{5cm}|p{7cm}|}
		\hline
		\textbf{Risk} & \textbf{Mitigation Strategy} \\
		\hline
		Requirement changes & Review project scope before implementation \\
		\hline
		Schedule delays & Conduct weekly progress monitoring \\
		\hline
		Merge conflicts & Apply Git branching strategy and code review \\
		\hline
		Incomplete documentation & Assign documentation responsibilities early \\
		\hline
		Integration issues & Perform continuous testing and integration \\
		\hline
	\end{tabular}
\end{table}

Risk management helps reduce project uncertainty and improves the likelihood of successful project completion.

\section{Summary}

This chapter presented the project management approach used in developing the Horse Racing Tournament Management System.

The project scope, MVP definition, team organization, project planning activities, Agile methodology, Jira-based task management, and risk management strategies were discussed.

These management practices provide a structured foundation for the successful execution of the project and support effective collaboration among team members.
